\seccion{1. Cardinalidad de las consultas}
Las relaciones iniciales son:\par 
$R(A,B)=\{(1,x),(2,y),(2,z),(3,x),(9,a)\}$\par 
$S(B,C,D)=\{(x,0,3),(y,2,1),(y,3,3),(w,3,0),(y,4,2)\}$.

\subsection*{C01. $R \times \rho_{B\_S \leftarrow B}(S)$}
Para poder realizar el producto cartesiano en herramientas como RelaX sin que exista conflicto por atributos homónimos, se renombra el atributo $B$ de la relación $S$ como $B\_S$. 
\[R \times \rho_{B\_S \leftarrow B}(S)\]
El producto tiene $5\cdot5=25$ tuplas. Se combina cada fila de $R$ con cada fila de $S'$ (la relación renombrada).
\begin{center}\begin{tabular}{ccccc}\toprule\rowcolor{Fondo} $A$ & $B$ & $B\_S$ & $C$ & $D$\\\midrule
1 & x & x & 0 & 3\\
1 & x & y & 2 & 1\\
1 & x & y & 3 & 3\\
1 & x & w & 3 & 0\\
1 & x & y & 4 & 2\\
2 & y & x & 0 & 3\\
2 & y & y & 2 & 1\\
2 & y & y & 3 & 3\\
2 & y & w & 3 & 0\\
2 & y & y & 4 & 2\\
2 & z & x & 0 & 3\\
2 & z & y & 2 & 1\\
2 & z & y & 3 & 3\\
2 & z & w & 3 & 0\\
2 & z & y & 4 & 2\\
3 & x & x & 0 & 3\\
3 & x & y & 2 & 1\\
3 & x & y & 3 & 3\\
3 & x & w & 3 & 0\\
3 & x & y & 4 & 2\\
9 & a & x & 0 & 3\\
9 & a & y & 2 & 1\\
9 & a & y & 3 & 3\\
9 & a & w & 3 & 0\\
9 & a & y & 4 & 2\\\bottomrule\end{tabular}\end{center}
\textbf{Cardinalidad: 25 tuplas.} No se eliminan filas duplicadas porque el producto cartesiano genera todas las combinaciones posibles.

\subsection*{C06. $\rho_{C \leftarrow A}(R) \bowtie S$}
Primero, el operador de renombramiento $\rho_{C \leftarrow A}(R)$ cambia el esquema de la relación $R$ de $(A,B)$ a $(C,B)$, manteniendo sus valores originales. La relación intermedia $R'$ queda como $R'(C,B)=\{(1,x),(2,y),(2,z),(3,x),(9,a)\}$.
\[ \rho_{C \leftarrow A}(R) \bowtie S \]
Al realizar la junta natural ($\bowtie$), se buscan coincidencias en \textbf{todos} los atributos comunes. En este caso, los atributos en común entre $R'$ y $S$ son dos: $B$ y $C$. Por lo tanto, se buscan tuplas donde tanto $B$ como $C$ sean idénticos.
\begin{center}\begin{tabular}{ccc}\toprule\rowcolor{Fondo} $B$ & $C$ & $D$\\\midrule
y & 2 & 1\\\bottomrule\end{tabular}\end{center}
\textbf{Cardinalidad: 1 tupla.} La única tupla de $R'$ que coincide simultáneamente con los valores de $B$ y $C$ de alguna tupla en $S$ es $(C=2, B=y)$, la cual empata con la tupla $(B=y, C=2, D=1)$ de $S$. Las demás filas no tienen coincidencia doble y se descartan.